\documentclass[12pt,a4paper]{article}
\usepackage[utf8]{inputenc}
\usepackage[T1]{fontenc}
\usepackage{amsmath,amssymb,amsthm}
\usepackage{geometry}
\usepackage{graphicx}
\usepackage{longtable}
\usepackage{array}
\usepackage{booktabs}
\usepackage{fancyhdr}
\usepackage{titlesec}
\usepackage{mathrsfs}

\geometry{margin=1in}

\pagestyle{fancy}
\fancyhf{}
\rhead{Cayley, Collected Papers, Vol.~XIII}
\lhead{Pages 201--250}
\cfoot{\thepage}

\titleformat{\section}{\large\bfseries}{\thesection.}{0.5em}{}
\titleformat{\subsection}{\normalsize\bfseries}{\thesubsection.}{0.5em}{}

\begin{document}

\begin{center}
{\Large\bfseries The Collected Mathematical Papers of\\[4pt]
Arthur Cayley, Vol.~XIII}\\[12pt]
{\large Pages 201--250}\\[8pt]
\rule{\textwidth}{0.4pt}
\end{center}

\tableofcontents
\newpage



\section*{921. On some Problems of Orthomorphosis.}
\addcontentsline{toc}{section}{\numberline{}921. On some Problems of Orthomorphosis.}

[From the \textit{Philosophical Transactions}, vol.~CLXXXIII (1892), pp.~653--669.]

\bigskip

\textbf{14.} To any line in figure 4, parallel to the axis $Ox$ or to the axis $Oy$, there
corresponds in figure 5 a bicircular quartic of the form $x^2 + y^2 - 2Ax^2 - 2By^2 - 1 = 0$.
The investigation is substantially the same as that contained in No.~12, and need
not be here given. But it is remarkable that also, to any line of figure 4 parallel
to a side of the square (that is, to any line $x \pm y = c$), there corresponds in figure 5
a bicircular quartic of the like form (for the sides of the square, or lines $x \pm y = \pm \sqrt{2}$
of figure 4, this bicircular quartic becomes the twice repeated circle $x^2 + y^2 - 1 = 0$
of figure 5, which is the result just obtained in No.~13). I investigate this result
as follows. Writing, as in No.~13,
\[
\operatorname{sn} x = p, \quad \operatorname{sn} iy = i\, \operatorname{sn} y = iq,
\]
we have, as above,
\[
x = \frac{p\sqrt{1-q^4} + q\sqrt{1-p^4}}{1-p^2q^2},
\quad
y = \frac{p\sqrt{1-q^4} - q\sqrt{1-p^4}}{1-p^2q^2},
\]
and thence
\[
x_1 = \frac{p\sqrt{1-q^4} + q\sqrt{1-p^4}}{1-p^2q^2},
\]
where $x_1 = x + y$ and $y_1 = x - y$.

Now assuming between $x$ and $y$ the relation $x + y = C$, and writing $\operatorname{sn} C = c$, this gives
\[
x_1 = \frac{p\sqrt{1-q^4} + q\sqrt{1-p^4}}{1-p^2q^2}.
\]
To obtain the required curve, we must between these equations (three independent
equations) eliminate $p$ and $q$. We have
\[
x_1 + y_1 = \frac{2p\sqrt{1-q^4}}{1-p^2q^2},
\quad
x_1 - y_1 = \frac{2q\sqrt{1-p^4}}{1-p^2q^2},
\]
and consequently, writing for convenience $n = 1 - p^2q^2$,
\[
nx_1 = p\sqrt{1-q^4}, \quad ny_1 = q\sqrt{1-p^4};
\]
these equations may be written
\[
p^2 = \frac{x_1^2}{n^2 + x_1^2}, \quad q^2 = \frac{y_1^2}{n^2 + y_1^2},
\]
and from these equations obtaining the expressions for $p^2$ and $q^2$, and thence the
expression for $p^2q^2 = 1 - n$, we find, after some easy reductions,
\[
n^2 - n(x_1^2 + y_1^2 - 1) - x_1^2y_1^2 = 0.
\]
But we have
\[
x_1^2 + y_1^2 = \frac{2(p^2 + q^2 - 2p^2q^2)}{(1-p^2q^2)^2},
\]
or substituting this value, the equation becomes
\[
\bigl[(x_1^2 + y_1^2)^2 - 2c^2(x_1^2 + y_1^2) + c^4\bigr] + (1-c^2)\bigl[(x_1^2 + y_1^2)^2 - c^2(x_1^2 + y_1^2)\bigr] - 4c^2x_1^2y_1^2 = 0,
\]
that is,
\[
(x_1^2 + y_1^2 - 1)^2 + (1-c^2)(x_1^2 + y_1^2 - c^2)(x_1^2 - y_1^2) = 0;
\]
the required equation. Transforming through an angle of $45^\circ$ by writing
\[
x_2 = \frac{x_1 + y_1}{\sqrt{2}}, \quad y_2 = \frac{x_1 - y_1}{\sqrt{2}}
\]
(where observe that the axis of $x_2$ is the line $FH$ of figure 5), the equation becomes
\[
(x_2^2 + y_2^2 - 1)^2 + (1-c^2)(2x_2^2 - 2y_2^2) = 0,
\]
or writing $c = \cos\gamma$ and therefore $1 - c^2 = \sin^2\gamma$, this equation becomes
\[
(x_2^2 + y_2^2)^2 - 2(x_2^2\cos^2\gamma + y_2^2\sin^2\gamma) + 1 = 0,
\]
a curve consisting of two indented ovals situate symmetrically in regard to the axis
$Fy_2$ of figure 5. In fact, writing in the equation $y_2 = 0$, we have for $x_2^2$ two real
positive values; but, writing $x_2 = 0$, we have for $y_2^2$ two imaginary values. For $\gamma=0$,
the equation becomes
\[
(x_2^2 + y_2^2)^2 - 2(x_2^2 + y_2^2) + 1 = 0,
\]
that is, we have the circle $x_2^2 + y_2^2 - 1 = 0$ twice repeated. One of the ovals is
shown in figure 5; the portion of it lying within the circle agrees with Schwarz's
figure, p.~113, turning this round through an angle of $45^\circ$.

For the lines $x - y = C$ of figure 4, we have in figure 5 the same system of
bicircular quartics turned round through an angle of $90^\circ$.

\bigskip
\textbf{II.}

\textbf{15.} I consider the general problem of the orthomorphosis of a circle into a
circle: we can, for the transformation of the circumference of the circle $x^2 + y^2 - 1 = 0$
into that of the circle $x_1^2 + y_1^2 - 1 = 0$, find a formula involving an arbitrary function
or (what is the same thing) an indefinite number of arbitrary constants. In fact,
writing for shortness $z = x + iy$, $z_1 = x_1 + iy_1$, and $\bar{z}$, $\bar{z}_1$ for the conjugate functions
$x - iy$, $x_1 - iy_1$; also $\phi(z)$ for a function of $z$ involving in general imaginary coefficients
$a + ib$, \&c., and $\bar{\phi}(\bar{z})$ for the like function with the conjugate coefficients
$a - ib$, \&c.; then if we assume
\[
z_1 = \frac{\phi(z)}{\bar{\phi}(\bar{z})} \cdot \frac{1}{z^m},
\]
where $m$ is any positive or negative integer, this implies
\[
\bar{z}_1 = \frac{\bar{\phi}(\bar{z})}{\phi(z)} \cdot \frac{1}{\bar{z}^m};
\]
consequently, if $x^2 + y^2 - 1 = 0$, that is, $z\bar{z} = 1$, or $\bar{z} = \frac{1}{z}$, we have
\[
z_1\bar{z}_1 = \frac{1}{z^m\bar{z}^m} = 1,
\]
or $z_1\bar{z}_1 = 1$, that is, $x_1^2 + y_1^2 - 1 = 0$. \hfill [(1)]

In a slightly different form, taking
\[
\phi(z) = (z - \alpha)(z - \beta)\ldots,
\]
$\alpha$, $\beta$, \&c., to denote any imaginary quantities,
and $\bar{\alpha}$, $\bar{\beta}$, \ldots\ the conjugate quantities; assuming
\[
\phi(z) = (z - \alpha)(z - \beta) \ldots,
\]
and taking $m$ for the number of factors, we have
\[
z_1 = \frac{(z - \alpha)(z - \beta)\ldots}{(\bar{z} - \bar{\alpha})(\bar{z} - \bar{\beta})\ldots} \cdot \frac{1}{z^m},
\]
and then (repeating the demonstration) we have
\[
z_1 = \frac{(z - \alpha)(z - \beta)\ldots}{(\frac{1}{z} - \bar{\alpha})(\frac{1}{z} - \bar{\beta})\ldots} \cdot \frac{1}{z^m},
\]
which, writing therein $\bar{z} = \frac{1}{z}$, becomes
\[
z_1 = \frac{(z - \alpha)(z - \beta)\ldots}{(1 - \bar{\alpha}z)(1 - \bar{\beta}z)\ldots},
\]
and consequently, if $z\bar{z} = 1$, then also $z_1\bar{z}_1 = 1$ as before.

We may in the expression for $z_1$ introduce a factor $\frac{A}{\bar{A}}$, or, what is the same
thing, a factor $A$ which is such that $A\bar{A} = 1$. In particular, we thus have the solution
\[
z_1 = A\frac{z - \alpha}{1 - \bar{\alpha}z}.
\]

\newpage
\textbf{16.} This is a solution with three arbitrary constants, viz.~$A$ (which may be
put $= \cos\lambda + i\sin\lambda$) is a single arbitrary constant, and $\alpha = a + ib$, is two arbitrary
constants; and these constants may be so determined that a given point in the
interior of the one circle, and a given point on the circumference thereof, shall
correspond respectively to a given point in the interior of the other circle and to a
given point on the circumference thereof. According to a well-known theorem of
Riemann's, any two simply connected areas included within given closed curves
respectively may be made to correspond to each other, and that in one way only,
under the foregoing conditions as to a pair of interior points and a pair of boundary
points: and we have, in what just precedes, the solution of the problem in the case
of two equal circles.

\textbf{17.} In the case of any other solution, we thus know that the correspondence
between the two circumferences cannot and does not imply a $(1, 1)$ correspondence
between the areas of the two circles: but it is interesting to enquire what happens.
I take a very particular case
\[
z_1 = \frac{z(z-2)}{1-2z},
\]
and therefore
\[
\bar{z}_1 = \frac{\bar{z}(\bar{z}-2)}{1-2\bar{z}},
\]
and consequently
\[
z_1\bar{z}_1 = \frac{z\bar{z}(z-2)(\bar{z}-2)}{(1-2z)(1-2\bar{z})},
\]
that is,
\[
x_1^2 + y_1^2 = \frac{(x^2+y^2)\{(x-2)^2+y^2\}}{(1-2x)^2+4y^2},
\]
so that, writing $x^2 + y^2 - 1 = 0$, we have $x_1^2 + y_1^2 - 1 = 0$, a correspondence of the two
circumferences. But to the circumference $x^2 + y^2 - 1 = 0$ there corresponds not only
the circumference $x_1^2 + y_1^2 - 1 = 0$, but another circumference. In fact, writing $x_1^2 + y_1^2 = 1$,
we have
\[
(x^2+y^2)\{(x-2)^2+y^2\} = (1-2x)^2 + 4y^2,
\]
that is,
\[
(x^2+y^2)^2 - 4x(x^2+y^2) + 4x - 1 = 0,
\]
or
\[
(x^2+y^2-1)(x^2+y^2-4x+1) = 0,
\]
and there is thus the other circle
\[
x^2+y^2-4x+1 = 0, \quad \text{or say} \quad (x-2)^2 + y^2 = 3,
\]
viz.~this is a circle, coordinates of centre $(2, 0)$ and radius $= \sqrt{3}$, cutting the circle
$x^2 + y^2 - 1 = 0$ in two real points. Referring to the figures 6 ($xy$) and 7 ($x_1y_1$), and
observing that $x_1 = 0$, $y_1 = 0$, that is, $z_1 = 0$, gives $z = 0$, or $z = 2$, that is, the points
$x = 0$, $y = 0$, and $x = 2$, $y = 0$, we see that to the centre $O$ in figure 7, there
correspond in figure 6 the points $O$, $M$ which are the centres of the two circles.

To any small closed curve, or say any small circle surrounding the point $O$ of
figure 7, there correspond in figure 6 small closed curves surrounding the points
$O$, $M$ respectively; and if in figure 7 the radius of the circle continually increases
and becomes nearly equal to unity, the closed curves of figure 6 continually increase,
changing at the same time their forms, and assume the forms shown by the dotted
lines of figure 6. It thus appears that, to the whole area of the circle $x^2 + y^2 - 1 = 0$
of figure 7, there correspond the two lunes $ACE$ and $ABD$ of figure 6; or if we
attend only to the area included within the circle $x^2 + y^2 - 1 = 0$ of this figure, then
there corresponds not the whole area of this circle, but only the area of the lune
$ACB$: and thus that the assumed relation $z_1 = \frac{z(z-2)}{1-2z}$ establishes, in fact, an ortho-\allowbreak
morphosis of the circle $x^2 + y^2 - 1 = 0$ into the lune $ACB$ which lies inside the circle
$x^2 + y^2 - 1 = 0$ and outside the circle $(x - 2)^2 + y^2 = 3 = 0$. It may be added that, to
the infinite area outside the circle $x^2 + y^2 - 1 = 0$ of figure 7, there correspond in
figure 6 first the area of the lens $AB$ common to the two circles, and secondly the
area outside the two circles: we have thus an orthomorphosis of the area outside
the circle $x_1^2 + y_1^2 - 1 = 0$ into these two areas respectively.

A somewhat more elegant example would have been that of the correspondence
\[
z_1 = \frac{z(z-\sqrt{2})}{1-\sqrt{2}\,z};
\]
here, corresponding to the circumference $x^2 + y^2 - 1 = 0$, we have the two equal
circumferences $x^2 + y^2 - 1 = 0$, and $(x - \sqrt{2})^2 + y^2 - 1 = 0$: and to the whole area of the
circle $x_1^2 + y_1^2 - 1 = 0$, there correspond two equal lunes $ACB$ and $ABD$.



\newpage
\section*{922. Note on the Lunar Theory.}
\addcontentsline{toc}{section}{\numberline{}922. Note on the Lunar Theory.}

[From the \textit{Monthly Notices of the Royal Astronomical Society}, vol.~LII. (1892), pp.~2--5.]

\bigskip

In the Lunar Theory, in whatever way worked out, the values ultimately obtained
for the coordinates $r$, $v$, $y$ should of course satisfy identically the equations of
motion; and that they do so, is the ultimate verification of the correctness of the
results obtained. It can hardly be hoped for that such a verification will ever be
made for Delaunay's results; and yet it would seem generally that the labour of
such a verification of the results to any extent, while exceeding (and possibly greatly
exceeding) that of obtaining these results by any method employed for that purpose,
ought still to be, so to speak, a labour of the same order. And one can, moreover,
imagine the process of verification having gone on pari passu with the obtainment of
the results.

The verification required is, in fact, that the values of $r$, $v$, $y$, obtained in any
manner, do satisfy identically the equations of motion; these are
\[
\frac{d^2x}{dt^2} + \frac{\mu x}{r^3} = \frac{\partial\Omega}{\partial x}, \quad
\frac{d^2y}{dt^2} + \frac{\mu y}{r^3} = \frac{\partial\Omega}{\partial y}, \quad
\frac{d^2z}{dt^2} + \frac{\mu z}{r^3} = \frac{\partial\Omega}{\partial z},
\]
where $\Omega$ is the disturbing function due to the Sun. I say the verification required is
that these equations are identically satisfied. This does not require the expressions
for $x$, $y$, $z$ to be given each in a single series, but only that the forms be such as
will allow of their being differentiated; and the comparison of the coefficients of the
cosines and sines of the same arguments on the two sides of each equation gives
the relations between the constants and the coefficients, which are in fact the
equations for the determination of these in the method of the variation of the
elements.

\newpage
\textbf{2.} As regards the coordinate $y$, the process of verification was, in fact,
carried out in the memoir \textit{On the Secular Acceleration of the Moon's Mean Motion},
\textit{Philosophical Transactions}, 1859, [516]; the result being that of the two equations
\[
\frac{dy}{dt} = \frac{\partial R}{\partial c^\prime}, \quad \frac{dc^\prime}{dt} + \frac{\partial R}{\partial y} = 0,
\]
the second is satisfied identically by the expressions for $y$ and $c^\prime$, while the first gives
$\frac{\partial R}{\partial c^\prime}$ equal to the calculated value of $\frac{dy}{dt}$.

\textbf{3.} We have $r = a(1 + \rho)$, and we thence obtain $r$ as a function of the time, the
form being in Delaunay's theory $r =$ function of $t + h + \varpi + \zeta$, where $\zeta$ denotes a sum of
cosines with constant coefficients: and thence, by integration, $v = n t + \epsilon +$ function
of the same arguments. The verification of the first two equations would thus
consist in showing that the expressions for $r$ and $v$ satisfy these equations; and the
process would seem to be one which, for any terms taken at pleasure, could be
carried out without too great labour.

\textbf{4.} But it would be interesting to see the process carried out, not for terms
taken at pleasure, but for the whole series as they present themselves in the actual
expressions for the coordinates. I observe that the verification for the coordinates
$r$, $v$ can be made without the explicit expression of $v$ in terms of the time; and that
the verification can be made to depend upon the identical equation $r^2\dot{v} = h$, or say
$\dot{v} = \dfrac{h}{r^2}$. We in fact have $r = \dfrac{a(1-e^2)}{1+e\cos v}$, and thence $r$ as a function of $v$, and we
have in like manner $\rho$, $v$ as functions of $t + h + \varpi + \zeta$; hence $r$, $\dot{v}$, functions of
$t + h + \varpi + \zeta$. The value of $\dot{v}$ obtained in this manner is equal to $\dfrac{h}{r^2}$, and this is
the required verification.

\newpage
\textbf{5.} To complete the verification for the coordinates $r$, $v$, we substitute the
value of $\dot{v}$, and also the value of $r$, in the first equation of motion, say
\[
\ddot{r} - r\dot{v}^2 + \frac{\mu}{r^2} = P;
\]
and the required result is that this equation shall be identically satisfied. The
expression for $r$ is in effect $r = a(1 + \rho)$, where $\rho$ is a sum of cosines; hence $\ddot{r}$ is a sum
of cosines, and the expression for $r\dot{v}^2$ is in like manner a sum of cosines: hence the
left-hand side is a sum of cosines; and the right-hand side $P$, inasmuch as $P$ is a
function of $r$, $v$ given as a sum of cosines, is also a sum of cosines. And if the
verification is to be made, it would consist in the comparison of the two sets of
cosines.

\textbf{6.} The verification should thus ultimately depend on the trigonometrical
identities, and the verification of them would in many cases be an interesting and
not too laborious exercise. As a very simple instance, we may take the following:
in the first approximation we have
\[
r = a(1 - e\cos nt), \quad v = nt + 2e\sin nt,
\]
and thence
\[
\dot{r} = ane\sin nt, \quad \dot{v} = n(1 + 2e\cos nt),
\]
and thus the equations become
\[
ane^2\cos nt - an^2(1 - e\cos nt)(1 + 4e\cos nt) + \frac{\mu}{a^2}(1 + 2e\cos nt) = P;
\]
or, when $e = 0$,
\[
-an^2 + \frac{\mu}{a^2} = 0,
\]
which is right; and, for the terms in $e$,
\[
ae^2n^2\cos nt + 3ane^2n^2\cos nt + \frac{\mu}{a^2}\cdot 2e\cos nt = P;
\]
and thus the equations become
\[
\frac{d^2r}{dt^2} - r\left(\frac{dv}{dt}\right)^2 + \frac{\mu}{r^2} = P, \quad
\frac{1}{r}\frac{d}{dt}\left(r^2\frac{dv}{dt}\right) = Q,
\]
where
\[
P = \frac{m^\prime}{a^{\prime 2}}\left(\frac{r}{a^{\prime}} - \frac{a^{\prime 2}r}{r^{\prime 3}}\cos(v - v^\prime)\right), \quad
Q = \frac{m^\prime}{a^{\prime 2}}\left(-\frac{a^{\prime 2}r}{r^{\prime 3}}\sin(v - v^\prime)\right),
\]
and the verification consists in showing that these equations are satisfied identically
by the assumed values of $r$ and $v$.



\newpage
\section*{923. Note on a Hyperdeterminant Identity.}
\addcontentsline{toc}{section}{\numberline{}923. Note on a Hyperdeterminant Identity.}

[From the \textit{Messenger of Mathematics}, vol.~xxi. (1892), pp.~131, 132.]

\bigskip

The following is in effect a well-known theorem; but I am not sure whether
it has been explicitly stated:

If $f$, $g$, $h$ are any functions of $x$, $y$, and we form the matrix
\[
\begin{pmatrix}
f & g & h\\
f_1 & g_1 & h_1\\
f_2 & g_2 & h_2
\end{pmatrix},
\]
where the suffixes 1 and 2 denote differentiation in regard to $x$ and $y$ respectively,
then the determinant
\[
\begin{vmatrix}
f & g & h\\
f_1 & g_1 & h_1\\
f_2 & g_2 & h_2
\end{vmatrix}
= 0
\]
is satisfied identically when $f$, $g$, $h$ are any linear functions of $x$, $y$; of course it
remains $=0$ when, instead of $x$, $y$, we write any functions whatever of $x$, $y$.

But if $f$, $g$, $h$ are any quadric functions of $x$, $y$, then writing $x$, $y = \dfrac{\xi}{\zeta}$,
$\dfrac{\eta}{\zeta}$, that is, considering $f$, $g$, $h$ as homogeneous functions of $(x, y, z)$, the
condition is
\[
\begin{vmatrix}
f & g & h\\
f_1 & g_1 & h_1\\
f_2 & g_2 & h_2
\end{vmatrix} = 0,
\]
which (as is obvious, and may be directly verified) is satisfied identically; but here
writing instead of $(x, y, z)$ any functions $(\xi, \eta, \zeta)$ of $(x, y, z)$, the condition no longer
remains satisfied identically; it gives in fact a differential equation of the second
order which must be satisfied by $\xi$, $\eta$, $\zeta$, in order that $f$, $g$, $h$, considered as functions
of $\xi$, $\eta$, $\zeta$, may be linear functions of these variables.

\newpage
We may in like manner consider $f$, $g$, $h$ as cubic functions of $(x, y)$; here, regarding
them as homogeneous functions of $(x, y, z)$, the conditions that they may be reducible
to linear functions of $(\xi, \eta, \zeta)$ are
\[
\begin{vmatrix}
f & g & h\\
f_1 & g_1 & h_1\\
f_2 & g_2 & h_2
\end{vmatrix} = 0,
\quad
\begin{vmatrix}
f & g & h\\
f_1 & g_1 & h_1\\
f_3 & g_3 & h_3
\end{vmatrix} = 0,
\quad
\begin{vmatrix}
f & g & h\\
f_2 & g_2 & h_2\\
f_3 & g_3 & h_3
\end{vmatrix} = 0,
\]
where 1, 2, 3 denote differentiation in regard to $x$, $y$, $z$ respectively. These are all
the same condition, and they are satisfied identically. But if we write for $(x, y, z)$
any functions $(\xi, \eta, \zeta)$ of $(x, y, z)$, then the conditions are differential equations which
must be satisfied by $\xi$, $\eta$, $\zeta$ in order that $f$, $g$, $h$, considered as functions of $\xi$, $\eta$, $\zeta$,
may be linear functions of these variables.



\newpage
\section*{924. On the Non-Existence of a Special Group of Points.}
\addcontentsline{toc}{section}{\numberline{}924. On the Non-Existence of a Special Group of Points.}

[From the \textit{Messenger of Mathematics}, vol.~xxi. (1892), pp.~132, 133.]

\bigskip

It is well known that, taking in a plane any eight points, every cubic through
these passes through a determinate ninth point. It is interesting to show that there
is no system of seven points such that every cubic through these passes through a
determinate eighth point.

Assuming such a system: first, no three of the points can be in a line: for, if
they were, then among the cubics through the seven points we have the line
through the three points and an arbitrary conic through the remaining four points,
and these composite cubics have no common eighth point of intersection.

Secondly, no six of the points can be on a conic: for, if they were, then among
the cubics through the seven points we have the conic through the six points and
an arbitrary line through the remaining point, and these composite cubics have no
common eighth point of intersection.

Taking now the points to be $1$, $2$, $3$, $4$, $5$, $6$, $7$; among the cubics through
these, we have the composite cubics $(A, P)$, $(B, Q)$, $(C, R)$, where $A$, $B$, $C$ are the
lines $67$, $75$, $56$, and $P$, $Q$, $R$ the conics $12345$, $12346$, $12347$ respectively; by what
precedes, the points $5$, $6$, $7$ do not lie on a line, and the points $(6, 7)$, $(7, 5)$ and
$(5, 6)$ neither of them lie on the conics $P$, $Q$, $R$ respectively.

The common eighth point of intersection, if it exists, must be
\[
(A, B, C);\quad (A, Q, R),\ (B, R, P),\ (C, P, Q);\quad
(P, B, C),\ (Q, C, A),\ (R, A, B);\quad \text{or}\ (P, Q, R).
\]
There is no point $(A, B, C)$.

There is no point $(A, Q, R)$: for $Q$, $R$ intersect only in the points $1$, $2$, $3$, $4$,
no one of which lies on $A$; and similarly, there is no point $(B, R, P)$ or $(C, P, Q)$.

$B$, $C$ intersect in $5$, which is a point on $P$; and thus $5$ is the only point
$(P, B, C)$. Similarly, $6$ is the only point $(Q, C, A)$ and $7$ is the only point $(R, A, B)$.

$P$, $Q$ intersect in $1$, $2$, $3$, $4$, which are each of them on $R$; hence $1$, $2$, $3$, $4$
are the only points $(P, Q, R)$.

Hence the points $1$, $2$, $3$, $4$, $5$, $6$, $7$ present themselves each once, and only
once, among the intersections of the three cubics, and there is no common eighth point
of intersection.



\newpage
\section*{925. On Waring's Formula for the Sum of the $m$th Powers of the Roots of an Equation.}
\addcontentsline{toc}{section}{\numberline{}925. On Waring's Formula for the Sum of the $m$th Powers of the Roots of an Equation.}

[From the \textit{Messenger of Mathematics}, vol.~xxi. (1892), pp.~133--137.]

\bigskip

The formula in question, Prob.~I.~of Waring's \textit{Meditationes Algebraic\ae}, Cambridge,
1782, making therein a slight change of notation, is as follows: viz. the equation being
\[
x^n - bx^{n-1} + cx^{n-2} - dx^{n-3} + \ldots = 0,
\]
then we have
\begin{align*}
(-)^{m-1} s_m &= mb^m - m\,c\\
&\quad - m\,b^{m-2}c + m(m-3)\,bd\\
&\quad + m\,b^{m-3}d - m(m-4)\,ce - m(m-4)\,bd + \tfrac{1}{2}m(m-4)(m-5)\,c^2d\\
&\quad + m\,b^{m-4}e + m(m-5)\,bf + m(m-5)\,ce\\
&\quad - m(m-5)(m-6)\,cd^2 - m(m-5)(m-6)\,be\\
&\quad + \tfrac{1}{6}m(m-5)(m-6)(m-7)\,c^3 - m\,g\\
&\quad - m(m-6)\,ch - m(m-6)\,dg - m(m-6)\,bf\\
&\quad + \tfrac{1}{2}m(m-6)(m-7)\,cd^2 - m(m-6)\,eg\\
&\quad + \tfrac{1}{2}m(m-6)(m-7)\,d^2f - \tfrac{1}{2}m(m-6)(m-7)\,ce^2\\
&\quad - \tfrac{1}{6}m(m-6)(m-7)(m-8)\,c^2e + \tfrac{1}{2}m(m-6)(m-7)\,d^3\\
&\quad + \tfrac{1}{24}m(m-6)(m-7)(m-8)(m-9)\,c^4 + \&\text{c.},
\end{align*}
where, reckoning the weights of $b$, $c$, $d$, $e$, \ldots\ as $1$, $2$, $3$, $4$, \ldots, respectively, the several
terms are all the terms of the weight $m$, or (what is the same thing) in the
coefficient of $b^{m-\theta}$ we have all the combinations of $c$, $d$, $e$, \ldots\ (or say all the
non-unitary combinations) of the weight $\theta$, and where the numerical coefficient of
\[
b^{m-\theta} c^d d^e e^f \ldots \quad (d + e + f + \ldots = \theta),
\]
is
\[
\frac{m \cdot m - (\theta + 1) \cdot m - (\theta + 2) \ldots m - (\theta - 1)}{c!\,d!\,e!\,f!\ldots}.
\]
Thus for the term $b^{m-8}c^2e^4$, $\theta = 8$; $c$, $d$, $e = 2$, $0$, $4$, $1$ respectively (the other exponents
each vanishing), and the coefficient is
\[
\frac{m \cdot m-6 \cdot m-7}{1 \cdot 2} = -\tfrac{1}{2}m(m-6)(m-7),
\]
as above; and so in other cases.

For the MacMahon form
\[
1 + bx + \frac{c}{1 \cdot 2}x^2 + \frac{d}{1 \cdot 2 \cdot 3}x^3 + \ldots = (1 - \alpha x)(1 - \beta x)(1 - \gamma x)\ldots,
\]
or say
\[
1 + bx + \frac{c}{1 \cdot 2}x^2 + \ldots = (1 - \alpha x)(1 - \beta x)\ldots,
\]
we must for $b$, $c$, $d$, \ldots\ write $b$, $\frac{1}{1 \cdot 2}c$, $\frac{1}{1 \cdot 2 \cdot 3}d$, \ldots\ respectively: we thus have
\begin{align*}
(-)^m \Pi(m-1) \cdot S_m &= \Pi(m-1) \cdot b^m\\
&\quad - \Pi(m-2) \cdot b^{m-2}c\\
&\quad + \Pi(m-3) \cdot b^{m-3}d\\
&\quad + \Pi(m-3) \cdot b^{m-3}cd\\
&\quad - \Pi(m-4) \cdot b^{m-4}e - \Pi(m-4) \cdot b^{m-4}c^2\\
&\quad + \&\text{c.},
\end{align*}
or say
\begin{multline*}
(-)^m \Pi(m-1) S_m = \Pi(m-1)b^m - \Pi(m-2)b^{m-2}c + \Pi(m-3)b^{m-3}d\\
+ \Pi(m-3)b^{m-3}cd - \Pi(m-4)b^{m-4}e - \Pi(m-4)b^{m-4}c^2 + \&\text{c.},
\end{multline*}
the numerical coefficient of
\[
b^{m-\theta} c^d d^e e^f \ldots \quad (d + e + f + \ldots = \theta)
\]
being
\[
\frac{(c + e + g + \ldots)m \cdot m - (\theta - \delta + 1) \cdot m - (\theta - \delta + 2) \ldots m - (\delta - 1)}{c!\,d!\,e!\ldots (1 \cdot 2)^c (1 \cdot 2 \cdot 3)^d (1 \cdot 2 \cdot 3 \cdot 4)^e \ldots}.
\]

\newpage
It is convenient to write down the literal terms in alphabetical order (AO),
calculating and affixing to each term the proper numerical coefficient; thus taking
\[
1 + bx + cx^2 + dx^3 + \ldots = (1 - \alpha x)(1 - \beta x)(1 - \gamma x)\ldots,
\]
we find
\begin{align*}
-120b^6 &= g \cdot 1\\
&\quad + f \cdot -6\\
&\quad + ce \cdot -15\\
&\quad + d^2 \cdot -10\\
&\quad + b^2e \cdot +30\\
&\quad + bcd \cdot +120\\
&\quad + c^3 \cdot +30\\
&\quad + b^3d \cdot -120\\
&\quad + b^2c^2 \cdot -270\\
&\quad + b^4c \cdot +360\\
&\quad + b^6 \cdot -120\\
&\quad \pm 541,
\end{align*}
this expression, as representing the value of the non-unitary function $S_6$, being in
fact a seminvariant.

It is to be remarked that the foregoing expression for the sum of the $m$th
powers of the roots of the equation
\[
x^n - bx^{n-1} + cx^{n-2} - \ldots = 0
\]
is, in fact, the series for $s_m$ continued so far only as the exponent of $b$ is not
negative: see as to this Note~XI.~of Lagrange's \textit{Equations Num\'eriques}. For the $a$
posteriori verification, observe that we have
\[
s_m = \alpha^m + \beta^m + \gamma^m + \ldots,
\]
or writing for a moment $u = -b$, say this is
\[
s_m = (-u)^m + (-u - c)^m + (-u - d)^m + \ldots,
\]
where
\[
f(x) = x^n - bx^{n-1} + cx^{n-2} - dx^{n-3} + \&\text{c.} = 0.
\]
Hence, by Lagrange's theorem,
\[
s_m = \frac{1}{2\pi i} \int \frac{x^m f^\prime(x)}{f(x)} dx,
\]
and the residue at infinity gives the required expression.



\newpage
\section*{926. Corrected Seminvariant Tables for the Weights 11 and 12.}
\addcontentsline{toc}{section}{\numberline{}926. Corrected Seminvariant Tables for the Weights 11 and 12.}

[From the \textit{Proceedings of the Royal Society}, vol.~L. (1891), pp.~431--437.]

\bigskip

The annexed Tables, extending those for the weights 1 to 10 given in my
``Memoir on Seminvariants,'' [875], require a correction. In the Tables, weight~10,
last two columns but one, instead of the columns $cdf - b^4d^2$, $ce^3 - c^5$, we ought to have
$ce^3 - c^5$, $cdf - b^4d^2$. The complete list up to the weight 12 is
\begin{align*}
w &= 1 && b\\
w &= 2 && c\\
w &= 3 && d,\ b^2\\
w &= 4 && e,\ bc\\
w &= 5 && f,\ bd,\ c^2\\
w &= 6 && g,\ be,\ cd,\ b^2c\\
w &= 7 && h,\ bf,\ ce,\ d^2,\ b^2d,\ bc^2\\
w &= 8 && i,\ bg,\ cf,\ de,\ b^2e,\ bcd,\ c^3,\ b^3c\\
w &= 9 && j,\ bh,\ cg,\ df,\ e^2,\ b^2f,\ bce,\ cd^2,\ bc^2d,\ b^3d,\ b^2c^2\\
w &= 10 && k,\ bi,\ ch,\ dg,\ ef,\ b^2g,\ bcf,\ bde,\ c^2d,\ b^3e,\\
&&& b^2cd,\ bc^3,\ b^4c;\quad cd^2f,\ ce^3 - c^5
\end{align*}

\newpage
\textbf{Seminvariant Tables, weight 11.}

\begin{longtable}{|c|c|c|c|c|c|c|c|c|c|c|c|c|c|c|}
\hline
$l$ & $c^2g$ & $cdef$ & $d^3$ & $b^2h$ & $bcf$ & $bce$ & $bd^2$ & $c^3d$ & $b^3i$ & $b^2cg$ & $b^2df$ & $b^2e^2$ & $bc^2e$ & $bcd^2$\\
\hline
\endhead
$bg$ & $0$ & $0$ & $0$ & $1$ & $0$ & $0$ & $0$ & $0$ & $0$ & $0$ & $0$ & $0$ & $0$ & $0$\\
$cf$ & $0$ & $0$ & $0$ & $0$ & $1$ & $0$ & $0$ & $0$ & $0$ & $0$ & $0$ & $0$ & $0$ & $0$\\
$de$ & $0$ & $0$ & $0$ & $0$ & $0$ & $1$ & $0$ & $0$ & $0$ & $0$ & $0$ & $0$ & $0$ & $0$\\
$c^3$ & $0$ & $0$ & $0$ & $0$ & $0$ & $0$ & $0$ & $1$ & $0$ & $0$ & $0$ & $0$ & $0$ & $0$\\
$b^2g$ & $0$ & $0$ & $0$ & $0$ & $0$ & $0$ & $0$ & $0$ & $1$ & $0$ & $0$ & $0$ & $0$ & $0$\\
$bcd$ & $0$ & $1$ & $0$ & $0$ & $0$ & $0$ & $0$ & $0$ & $0$ & $0$ & $0$ & $0$ & $0$ & $0$\\
$b^2f$ & $0$ & $0$ & $0$ & $0$ & $0$ & $0$ & $0$ & $0$ & $0$ & $0$ & $1$ & $0$ & $0$ & $0$\\
$b^2e$ & $0$ & $0$ & $0$ & $0$ & $0$ & $0$ & $0$ & $0$ & $0$ & $0$ & $0$ & $1$ & $0$ & $0$\\
$bc^2$ & $1$ & $0$ & $0$ & $0$ & $0$ & $0$ & $0$ & $0$ & $0$ & $0$ & $0$ & $0$ & $0$ & $0$\\
$b^3d$ & $0$ & $0$ & $0$ & $0$ & $0$ & $0$ & $0$ & $0$ & $0$ & $0$ & $0$ & $0$ & $0$ & $0$\\
$b^3c$ & $0$ & $0$ & $0$ & $0$ & $0$ & $0$ & $0$ & $0$ & $0$ & $0$ & $0$ & $0$ & $0$ & $0$\\
\hline
\end{longtable}

\newpage
\textbf{Seminvariant Tables, weight 12 (part 1).}

\begin{longtable}{|c|c|c|c|c|c|c|c|c|c|c|c|c|c|c|}
\hline
$m$ & $c^3f$ & $c^2e^2$ & $c^2d^2$ & $b^2i$ & $bch$ & $bcg$ & $bdf$ & $be^2$ & $c^4$ & $b^3j$ & $b^2ci$ & $b^2dh$ & $b^2eg$ & $b^2f^2$\\
\hline
\endhead
$bi$ & $0$ & $0$ & $0$ & $1$ & $0$ & $0$ & $0$ & $0$ & $0$ & $0$ & $0$ & $0$ & $0$ & $0$\\
$ch$ & $0$ & $0$ & $0$ & $0$ & $1$ & $0$ & $0$ & $0$ & $0$ & $0$ & $0$ & $0$ & $0$ & $0$\\
$dg$ & $0$ & $0$ & $0$ & $0$ & $0$ & $1$ & $0$ & $0$ & $0$ & $0$ & $0$ & $0$ & $0$ & $0$\\
$ef$ & $0$ & $0$ & $0$ & $0$ & $0$ & $0$ & $1$ & $0$ & $0$ & $0$ & $0$ & $0$ & $0$ & $0$\\
$c^2d$ & $1$ & $0$ & $0$ & $0$ & $0$ & $0$ & $0$ & $0$ & $0$ & $0$ & $0$ & $0$ & $0$ & $0$\\
$b^2h$ & $0$ & $0$ & $0$ & $0$ & $0$ & $0$ & $0$ & $0$ & $0$ & $1$ & $0$ & $0$ & $0$ & $0$\\
$bcg$ & $0$ & $0$ & $0$ & $0$ & $0$ & $0$ & $0$ & $0$ & $0$ & $0$ & $1$ & $0$ & $0$ & $0$\\
$bcf$ & $0$ & $0$ & $0$ & $0$ & $0$ & $0$ & $0$ & $0$ & $0$ & $0$ & $0$ & $1$ & $0$ & $0$\\
$bce$ & $0$ & $0$ & $0$ & $0$ & $0$ & $0$ & $0$ & $0$ & $0$ & $0$ & $0$ & $0$ & $1$ & $0$\\
$bd^2$ & $0$ & $0$ & $0$ & $0$ & $0$ & $0$ & $0$ & $0$ & $0$ & $0$ & $0$ & $0$ & $0$ & $0$\\
$c^3$ & $0$ & $1$ & $0$ & $0$ & $0$ & $0$ & $0$ & $0$ & $0$ & $0$ & $0$ & $0$ & $0$ & $0$\\
$b^3g$ & $0$ & $0$ & $0$ & $0$ & $0$ & $0$ & $0$ & $0$ & $0$ & $0$ & $0$ & $0$ & $0$ & $0$\\
$b^3f$ & $0$ & $0$ & $0$ & $0$ & $0$ & $0$ & $0$ & $0$ & $0$ & $0$ & $0$ & $0$ & $0$ & $0$\\
$b^3e$ & $0$ & $0$ & $0$ & $0$ & $0$ & $0$ & $0$ & $0$ & $0$ & $0$ & $0$ & $0$ & $0$ & $0$\\
$b^2c^2$ & $0$ & $0$ & $0$ & $0$ & $0$ & $0$ & $0$ & $0$ & $1$ & $0$ & $0$ & $0$ & $0$ & $0$\\
\hline
\end{longtable}

\newpage
\textbf{Seminvariant Tables, weight 12 (part 2).}

\begin{longtable}{|c|c|c|c|c|c|c|c|c|c|c|c|c|c|c|}
\hline
$bce$ & $bcd$ & $c^2d$ & $c^5$ & $b^4k$ & $b^3bi$ & $b^3ch$ & $b^3dg$ & $b^3ef$ & $b^2c^2g$ & $b^2cdf$ & $b^2ce^2$ & $b^2d^3$ & $bc^3f$ & $bc^2e^2$\\
\hline
\endhead
$b^2i$ & $0$ & $0$ & $0$ & $1$ & $0$ & $0$ & $0$ & $0$ & $0$ & $0$ & $0$ & $0$ & $0$ & $0$\\
$bcg$ & $0$ & $0$ & $0$ & $0$ & $0$ & $0$ & $0$ & $0$ & $1$ & $0$ & $0$ & $0$ & $0$ & $0$\\
$bce$ & $0$ & $0$ & $0$ & $0$ & $0$ & $0$ & $0$ & $0$ & $0$ & $0$ & $0$ & $0$ & $0$ & $0$\\
$bd^2$ & $0$ & $0$ & $0$ & $0$ & $0$ & $0$ & $0$ & $0$ & $0$ & $0$ & $0$ & $0$ & $0$ & $0$\\
$c^3$ & $0$ & $0$ & $0$ & $0$ & $0$ & $0$ & $0$ & $0$ & $0$ & $0$ & $0$ & $0$ & $0$ & $0$\\
$b^3g$ & $0$ & $0$ & $0$ & $0$ & $0$ & $0$ & $0$ & $0$ & $0$ & $0$ & $0$ & $0$ & $0$ & $0$\\
$b^3f$ & $0$ & $0$ & $0$ & $0$ & $0$ & $0$ & $0$ & $0$ & $0$ & $0$ & $0$ & $0$ & $0$ & $0$\\
$b^3e$ & $0$ & $0$ & $0$ & $0$ & $0$ & $0$ & $0$ & $0$ & $0$ & $0$ & $0$ & $0$ & $0$ & $0$\\
$b^2c^2$ & $0$ & $0$ & $0$ & $0$ & $0$ & $0$ & $0$ & $0$ & $0$ & $0$ & $0$ & $0$ & $0$ & $0$\\
\hline
\end{longtable}



\newpage
\section*{927. On Clifford's Paper ``On Syzygetic Relations among the Powers of Linear Quantics.''}
\addcontentsline{toc}{section}{\numberline{}927. On Clifford's Paper ``On Syzygetic Relations among the Powers of Linear Quantics.''}

[From the \textit{Proceedings of the Royal Society}, vol.~L. (1891), pp.~437--439.]

\bigskip

The question discussed in this paper is as follows: writing $U = (a, b, c, \ldots \mid x, y)^n$,
a binary quantic of the order $n$, with the coefficients $(a, b, c, \ldots)$; and taking $k$ a given
positive integer, to find the number of independent syzygies of the form
\[
\alpha U^k + \beta U^{k-1}V + \gamma U^{k-2}V^2 + \ldots = 0,
\]
where $\alpha$, $\beta$, $\gamma$, \ldots\ are rational and integral functions of the coefficients of the
quantics $U$ and $V$, and where the degree in these coefficients of each term $U^{k-f}V^f$
is the same for all values of $f$; the syzygy being thus an identical relation between
the different powers $U^k$, $U^{k-1}V$, \ldots\ of the quantics $U$, $V$.

\textbf{1.} I consider in particular the case where the quantics $U$, $V$ are linear quantics,
say $U = ax + by$, $V = cx + dy$. Here taking $k = 2$, we have between $U^2$, $UV$,
and $V^2$ a single relation
\[
(ad - bc)^2 U^2 \cdot V^2 = (ad - bc)^2 (UV)^2,
\]
that is,
\[
(ad - bc)^2 U^2 \cdot V^2 - (ad - bc)^2 (UV)^2 = 0,
\]
which may be written
\[
\alpha U^2 + \beta UV + \gamma V^2 = 0,
\]
where $\alpha = (ad - bc)^2 V^2$, $\beta = -2(ad - bc)^2 UV$, $\gamma = (ad - bc)^2 U^2$, are of the
proper degrees in the coefficients; viz.~considering the degrees in $(a, b)$ and in
$(c, d)$ respectively, the degrees of $\alpha$, $\beta$, $\gamma$ are as follows:
\begin{center}
\begin{tabular}{c|cc}
& $(a, b)$ & $(c, d)$ \\
\hline
$\alpha$ & $0$, & $2$ \\
$\beta$ & $1$, & $1$ \\
$\gamma$ & $2$, & $0$
\end{tabular}
\end{center}
and the degrees of $\alpha U^2$, $\beta UV$, $\gamma V^2$ are in $(a, b)$ each $= 2$, and in $(c, d)$
each $= 2$, as they should be.

\newpage
\textbf{2.} But we see further that the twofold relation
\[
\alpha U^2 + \beta UV + \gamma V^2 = 0, \quad \alpha^\prime U^2 + \beta^\prime UV + \gamma^\prime V^2 = 0,
\]
where the coefficients $\alpha^\prime$, $\beta^\prime$, $\gamma^\prime$ are independent of $\alpha$, $\beta$, $\gamma$, cannot be satisfied
for arbitrary values of $(a, b, c, d)$; for from the two equations eliminating $U^2$:
$UV$:
$V^2$, we should obtain
\[
\alpha\beta^\prime - \alpha^\prime\beta = 0, \quad \alpha\gamma^\prime - \alpha^\prime\gamma = 0, \quad \beta\gamma^\prime - \beta^\prime\gamma = 0,
\]
which are three relations between the coefficients $\alpha$, $\beta$, \ldots, $\gamma^\prime$.

Hence there is not in general any twofold syzygy; and the conclusion is that
for linear quantics, and for the powers $U^2$, $UV$, $V^2$, the only syzygy is the foregoing
single one.

\textbf{3.} I consider, more generally, the syzygies between the powers $U^k$,
$U^{k-1}V$, \ldots, $V^k$ of the two linear quantics $U$, $V$; or say the syzygies between the
$k + 1$ powers $(U, V)^k$. The number of terms is $k + 1$, and the number of
coefficients of the quantics is $4$; hence the number of coefficients of the products
$U^k$, $U^{k-1}V$, \ldots\ is $= (k + 1) + k = 2k + 1$; and we have thus $2k + 1$ products of
the coefficients, each of which must have the degree $k$ in $(a, b)$ and the degree $k$
in $(c, d)$. The number of products of the proper degrees is $= k + 1$, and this is the
number of independent syzygies.

The number of syzygies is thus $= k + 1$, and since the number of terms is
also $= k + 1$, the syzygies give rise to a single equation; that is, there is a single
relation
\[
A \cdot U^k + B \cdot U^{k-1}V + \ldots = 0,
\]
where the coefficients $A$, $B$, \ldots\ are products of the degrees just referred to, the
number of them being $= k + 1$, but they are connected by $k$ relations, so that there
remains only a single constant, that is, we have a single syzygy.

\textbf{4.} But the conclusion is different if we consider the syzygies between the
powers $U^k$, $U^{k-1}V$, \ldots\ of quantics $U$, $V$ of any order $n$. Considering the case
where $U$, $V$ are quadric quantics, say $U = (a, b, c \mid x, y)^2$, $V = (d, e, f \mid x, y)^2$,
we have between the $k + 1$ powers $U^k$, $U^{k-1}V$, \ldots, $V^k$, a number of syzygies which
is at least $= k + 1$, and may be greater. In fact, writing $\Theta = (ad - bf)^2 - 4(ae - bd)(cf - ae)$,
the condition for the two quantics to have a common factor is $\Theta = 0$; and when
this is so, we have $U = (\alpha x + \beta y)(a^\prime x + b^\prime y)$, $V = (\alpha x + \beta y)(c^\prime x + d^\prime y)$,
so that $V = \dfrac{c^\prime x + d^\prime y}{a^\prime x + b^\prime y} \cdot U$, or say $V = \dfrac{W}{W^\prime} \cdot U$; hence $VU^\prime - UV^\prime = 0$,
that is, the quantics have a common factor; and we thence obtain the like relations
between the powers of $U$ and $V$. But it is not in this way that we obtain the
syzygies in question; and the general theory requires a further examination.



\newpage
\section*{928. On the Analytical Theory of the Congruency.}
\addcontentsline{toc}{section}{\numberline{}928. On the Analytical Theory of the Congruency.}

[From the \textit{Proceedings of the Royal Society}, vol.~L. (1891), pp.~440--457.]

\bigskip

The theory of the congruency, or congruence, of lines was established by
Kummer in his celebrated memoir ``Ueber die algebraischen Strahlensysteme, in
Welchen die Strahlen durch eine algebraische Gleichung verbunden sind,'' \textit{Berlin.
Abh.} 1866. The congruency is a system of lines satisfying two conditions, or say
a twofold relation; and it thus occupies the same position in relation to the complex
(a system of lines satisfying one condition, or say a onefold relation) that a surface
does in relation to a curve.

\textbf{1.} I consider a line determined by its six coordinates $(a, b, c, f, g, h)$, where
$af + bg + ch = 0$; and I represent the consecutive line by $(a + \delta a, \ldots, f + \delta f, \ldots)$.
The condition in order that the two lines may intersect is
\[
(a + \lambda a^\prime,\ b + \lambda b^\prime,\ c + \lambda c^\prime,\ f + \lambda f^\prime,\ g + \lambda g^\prime,\ h + \lambda h^\prime),
\]
viz.~the consecutive line $(a_1, b_1, c_1, f_1, g_1, h_1)$ belongs to the linear congruency defined
by these two equations.

Forming with these a linear combination, the equation of the congruency may
be written
\[
\begin{vmatrix}
a & b & c & f & g & h \\
a^\prime & b^\prime & c^\prime & f^\prime & g^\prime & h^\prime \\
\delta a & \delta b & \delta c & \delta f & \delta g & \delta h
\end{vmatrix} = 0,
\]
which is the condition for the intersection of the line $(a, b, c, f, g, h)$ with the line
$(\delta a, \delta b, \delta c, \delta f, \delta g, \delta h)$.

\textbf{2.} For the point $2$, we have
\begin{align*}
hy - gz + aw &= 0, \quad h^\prime y - g^\prime z + a^\prime w = 0,\\
-hx\quad + fz + bw &= 0, \quad -h^\prime x\quad + f^\prime z + b^\prime w = 0,\\
gx - fy\quad\; + cw &= 0, \quad g^\prime x - f^\prime y\quad\; + c^\prime w = 0,
\end{align*}
where $(x, y, z, w)$ are the coordinates of the point. Hence we have
\[
\frac{x}{fw - bg} = \frac{y}{ah - cf} = \frac{z}{ab - ch} = \frac{w}{af + bg + ch},
\]
and similarly for the other point. The line of the congruency is thus determined
by the two points in which it meets the two planes.

\textbf{3.} The order of the congruency is the number of lines which pass through
a given point. Taking the point to be $(x_0, y_0, z_0, w_0)$, the condition that a line of the
congruency passes through this point is that the coordinates $(a, b, c, f, g, h)$ satisfy
the two equations
\[
af + bg + ch = 0, \quad hx_0 - gy_0 + fz_0 + aw_0 = 0,
\]
and the order is thus the number of solutions of these equations, which is equal
to the number of intersections of the two quadric surfaces, viz.~it is $= 4$.

\textbf{4.} The class of the congruency is the number of lines which lie in a given
plane. Taking the plane to be $lx + my + nz + pw = 0$, the condition that a line
of the congruency lies in this plane is that the coordinates $(a, b, c, f, g, h)$ satisfy
the two equations
\[
af + bg + ch = 0, \quad la + mb + nc + pf = 0,
\]
and the class is thus the number of solutions of these equations, which is equal
to the number of intersections of the two quadric surfaces, viz.~it is $= 4$.

\textbf{5.} The congruency has a singular surface, which is the locus of points
through which pass an infinite number of lines of the congruency. This surface
is of the order equal to the number of lines of the congruency which pass through
a given point, and which are contained in a given plane through the point; that is,
it is the number of intersections of the line of the congruency with the plane, which
is $= 4$.

\textbf{6.} The theory may be extended to the case of a congruency of higher
order, where the lines satisfy three or more conditions. The general theory of such
congruencies presents many interesting questions, which have not yet been fully
investigated.



\newpage
\section*{929. Note on the Skew Surfaces Applicable upon a Given Skew Surface.}
\addcontentsline{toc}{section}{\numberline{}929. Note on the Skew Surfaces Applicable upon a Given Skew Surface.}

[From the \textit{Philosophical Transactions}, vol.~CLXXXIII (1892), pp.~677--688.]

\bigskip

The question of the applicability of one skew surface upon another is a
problem of great interest in the theory of surfaces. The general solution depends
upon the integration of a partial differential equation of the second order; and the
particular case where the given surface is a skew surface, or ruled surface, presents
special features which make it worthy of separate consideration.

\textbf{1.} I consider a skew surface generated by the motion of a straight line;
and I seek to determine the skew surfaces which are applicable upon it. The
general theory of applicable surfaces shows that two surfaces are applicable when
their linear elements can be expressed in the same form; that is, when the
geodesic distances between corresponding points are equal.

For a skew surface, the linear element can be expressed in the form
\[
ds^2 = du^2 + 2\cos\theta\,du\,dv + (u^2 + \alpha u + \beta)\,dv^2,
\]
where $\theta$ is the angle between the generating line and the direction of the parameter
$v$, and $\alpha$, $\beta$ are functions of $v$ only. The condition for applicability is that this
expression for $ds^2$ shall be the same for the two surfaces, or at least reducible to
the same form.

\textbf{2.} If we consider the strip of the surface bounded by two consecutive
generating lines, the distance between these lines, measured along their common
perpendicular, is the parameter of distribution. The condition of applicability
then requires that the strips of the two surfaces shall have the same breadth (or
distance on each side from $P$) of the strip or element of the surface, the same
angle between the bounding lines, and the same parameter of distribution.

\textbf{3.} For the general skew surface, the condition of applicability leads to a
partial differential equation. Writing the equation of the surface in the form
\[
z = f(x, y),
\]
the condition is that the first and second fundamental quantities shall satisfy the
equation of Gauss. The integration of this equation gives the most general
surface applicable upon the given surface; and it appears that the solution depends
upon an arbitrary function.

\textbf{4.} In the particular case where the given surface is a surface of translation,
the problem is simplified. A surface of translation is generated by the translation
of a curve parallel to itself, the directing curve being also moved in such a manner
that the generating curve always remains parallel to a fixed plane. For such a
surface, the linear element can be expressed in a simple form, and the condition of
applicability can be more easily investigated.

\textbf{5.} The case where the given surface is a helicoid is also of special interest.
A helicoid is generated by the motion of a straight line which turns uniformly
round a fixed axis, while at the same time it moves uniformly in the direction of
that axis. The linear element of a helicoid can be expressed in the form
\[
ds^2 = dr^2 + r^2\,d\theta^2 + (c\,d\theta + dz)^2,
\]
and the condition of applicability upon another helicoid can be completely
determined. It appears that the only surfaces applicable upon a helicoid are
helicoids; and the relation between the two helicoids can be expressed in a very
simple form.

\textbf{6.} More generally, if the given skew surface is a surface of revolution,
the problem can be solved by the use of elliptic integrals. The linear element of
a surface of revolution can be expressed in the form
\[
ds^2 = dr^2 + r^2\,d\theta^2 + dz^2,
\]
where $z$ is a function of $r$; and the condition of applicability upon another surface
of revolution leads to an equation which can be integrated by elliptic functions.

\textbf{7.} The general theory of the applicability of skew surfaces is connected
with the theory of the deformation of surfaces, and with the theory of the
integration of partial differential equations. The problem presents many difficulties,
which have not yet been completely overcome; but the results which have been
obtained are sufficient to show the great interest and importance of the subject.



\newpage
\section*{930. Sur la Surface des Ondes.}
\addcontentsline{toc}{section}{\numberline{}930. Sur la Surface des Ondes.}

[From the \textit{Annali di Matematica}, Ser.~II., vol.~xx. (1892), pp.~1--10.]

\bigskip

\textbf{1.} Je me sers de la d\'etermination bien connue de Fresnel de la surface
des ondes, c'est-\`a-dire de l'enveloppe du plan
\[
lx + my + nz = 1,
\]
o\`u $l$, $m$, $n$ sont les cosinus directeurs d'un rayon de l'onde; le plan est donn\'e par
l'\'equation
\[
\frac{a^2l^2}{1-a^2r^2} + \frac{b^2m^2}{1-b^2r^2} + \frac{c^2n^2}{1-c^2r^2} = 0,
\]
o\`u $r^2 = l^2 + m^2 + n^2$; et l'\'equation de la surface s'obtient en \'eliminant $l$, $m$, $n$, $r$
entre ces deux \'equations et les \'equations qui expriment que le plan touche la surface.

\textbf{2.} On peut, au lieu de cela, \'ecrire
\[
lx + my + nz = 1, \quad l^2 + m^2 + n^2 = r^2,
\]
et consid\'erer $r$ comme un param\`etre. L'enveloppe s'obtient alors en \'eliminant
$r$ entre l'\'equation du plan et son d\'eriv\'ee par rapport \`a $r$.

\textbf{3.} En \'ecrivant
\[
\frac{l^2}{1-a^2r^2} + \frac{m^2}{1-b^2r^2} + \frac{n^2}{1-c^2r^2} = 0,
\]
on a l'\'equation de la surface sous la forme
\[
\frac{a^2x^2}{r^2-a^2} + \frac{b^2y^2}{r^2-b^2} + \frac{c^2z^2}{r^2-c^2} = 1,
\]
o\`u $r$ est d\'etermin\'e par l'\'equation
\[
\frac{x^2}{r^2-a^2} + \frac{y^2}{r^2-b^2} + \frac{z^2}{r^2-c^2} = 1.
\]

\textbf{4.} Pour obtenir l'\'equation de la surface en coordonn\'ees rectangulaires,
il faut \'eliminer $r$ entre ces deux \'equations. On obtient ainsi une \'equation du
degr\'e 16; mais cette \'equation se r\'eduit au degr\'e 4 lorsqu'on consid\`ere les
carr\'es des coordonn\'ees comme les variables.

\textbf{5.} Je remarque que l'\'equation de la surface peut s'\'ecrire sous la forme
\[
\begin{vmatrix}
x^2 & y^2 & z^2 \\
r^2-a^2 & r^2-b^2 & r^2-c^2 \\
1 & 1 & 1
\end{vmatrix} = 0,
\]
et que cette forme met en \'evidence plusieurs propri\'et\'es importantes de la surface.

\textbf{6.} On a, pour les coordonn\'ees du point de contact du plan tangent,
\[
\xi = \frac{a^2l}{1-a^2r^2}, \quad \eta = \frac{b^2m}{1-b^2r^2}, \quad \zeta = \frac{c^2n}{1-c^2r^2}.
\]
La premi\`ere de ces \'equations donne
\[
\xi(1-a^2r^2) = a^2l,
\]
et de l\`a
\[
l = \frac{\xi(1-a^2r^2)}{a^2}.
\]
On a aussi
\[
x = \frac{l}{r^2-a^2}, \quad y = \frac{m}{r^2-b^2}, \quad z = \frac{n}{r^2-c^2},
\]
et de l\`a
\[
u + v = (b+c)\lambda^2 + (c+a)\mu^2 + (a+b)\nu^2,
\]
o\`u $\lambda$, $\mu$, $\nu$ sont les cosinus directeurs de la normale.

\newpage
\textbf{7.} On a
\[
\eta - b = T(f-a)\{\xi\eta + \zeta - (\eta + \zeta)\xi\},
\]
et de l\`a
\[
\eta - b = T(f-a)(\xi\eta + \zeta - \eta\xi - \zeta\xi),
\]
o\`u $T$ est une fonction de $r$.

\textbf{8.} La surface des ondes a plusieurs propri\'et\'es remarquables. En
particulier, elle poss\`ede quatre points coniques, qui correspondent aux directions
des axes optiques; et elle a aussi quatre droites doubles, qui sont les intersections
de la surface avec le plan \`a l'infini.

\textbf{9.} Les courbes de contact du plan tangent sont des courbes du quatri\`eme
ordre, qui ont des points singuliers aux points coniques de la surface. Ces courbes
sont des courbes gauches, et leur \'etude pr\'esente un grand int\'er\^et.

\textbf{10.} La surface des centres de la surface des ondes est une surface du
vingt-quatri\`eme ordre. Elle a des points singuliers qui correspondent aux points
singuliers de la surface des ondes; et son \'etude est intimement li\'ee \`a celle de la
surface des ondes.

\textbf{11.} On peut \'etudier la surface des ondes au moyen de ses sections par
des plans parall\`eles aux coordonn\'ees. Ces sections sont des courbes du quatri\`eme
ordre, qui ont des propri\'et\'es remarquables. En particulier, les sections par des
plans parall\`eles aux plans de coordonn\'ees sont des courbes ayant des points doubles
aux points o\`u le plan coupe les axes.

\textbf{12.} L'\'equation diff\'erentielle des courbes de contact peut s'obtenir en
\'eliminant $r$ entre l'\'equation du plan et ses d\'eriv\'ees. On obtient ainsi une
\'equation du premier ordre, qui d\'etermine les courbes de contact sur la surface.

\textbf{13.} La surface des ondes peut aussi \^etre \'etudi\'ee au moyen de ses
courbures principales. Les rayons de courbure principaux sont donn\'es par une
\'equation du deuxi\`eme degr\'e, qui s'obtient en \'egalant \`a z\'ero le d\'eterminant des
deriv\'ees secondes de la fonction qui d\'efinit la surface.

\textbf{14.} Le point de contact du plan tangent peut \^etre d\'etermin\'e au moyen
des formules donn\'ees ci-dessus. On a
\[
X = \mu Z - \nu Y, \quad Y = \nu X - \lambda Z, \quad Z = \lambda Y - \mu X,
\]
et de l\`a
\[
X : Y : Z = \lambda : \mu : \nu,
\]
c'est-\`a-dire que le rayon vecteur du point de contact est parall\`ele \`a la normale au
plan tangent.

\newpage
\textbf{15.} De la valeur ci-dessus donn\'ee pour $\eta - b$, on d\'eduit que la surface
a des propri\'et\'es de sym\'etrie remarquables. En particulier, elle est sym\'etrique
par rapport aux plans de coordonn\'ees, et aussi par rapport \`a l'origine.

\textbf{16.} L'\'equation de la surface peut aussi s'\'ecrire sous la forme
\[
\frac{x^2}{a^2-r^2} + \frac{y^2}{b^2-r^2} + \frac{z^2}{c^2-r^2} = 1,
\]
o\`u $r$ est d\'etermin\'e par l'\'equation
\[
\frac{a^2x^2}{(a^2-r^2)^2} + \frac{b^2y^2}{(b^2-r^2)^2} + \frac{c^2z^2}{(c^2-r^2)^2} = 0.
\]
Cette forme de l'\'equation met en \'evidence la connexion entre la surface des ondes
et les surfaces confocales du second degr\'e.

\textbf{17.} On peut consid\'erer la surface des ondes comme une surface de la
classe 16, ayant 4 points coniques et 4 droites doubles. Le genre de la surface
est $= 3$, et ses invariants num\'eriques peuvent \^etre calcul\'es par les formules de la
g\'eom\'etrie \'enum\'erative.

\textbf{18.} La surface des ondes a \'et\'e l'objet de nombreuses recherches. Les
r\'esultats les plus importants sont dus \`a Fresnel, Hamilton, Kummer, et Cayley.
La th\'eorie de cette surface est intimement li\'ee \`a celle de la propagation de la
lumi\`ere dans les milieux cristallins.

\textbf{19.} On peut \'etudier la surface des ondes au moyen des coordonn\'ees
elliptiques. En posant
\[
\frac{x^2}{\theta-a^2} + \frac{y^2}{\theta-b^2} + \frac{z^2}{\theta-c^2} = 1,
\]
on d\'efinit trois syst\`emes de surfaces orthogonales; et les propri\'et\'es de la surface
des ondes peuvent \^etre d\'eduites de la th\'eorie g\'en\'erale de ces coordonn\'ees.

\textbf{20.} L'\'equation diff\'erentielle des courbes de contact peut se mettre sous
la forme
\[
A\,dx^2 + B\,dy^2 + C\,dz^2 + 2D\,dy\,dz + 2E\,dz\,dx + 2F\,dx\,dy = 0,
\]
o\`u les coefficients $A$, $B$, $C$, $D$, $E$, $F$ sont des fonctions de $x$, $y$, $z$. Cette \'equation
d\'etermine les deux directions principales en chaque point de la surface.

\newpage
\textbf{21.} On obtient sans peine, au moyen des valeurs trouv\'ees ci-dessus,
l'\'equation de la surface des ondes sous la forme explicite. L'\'equation finale est
\[
(a^2x^2 + b^2y^2 + c^2z^2)(x^2 + y^2 + z^2)^2 - (a^2 + b^2 + c^2)(x^2 + y^2 + z^2)(a^2x^2 + b^2y^2 + c^2z^2)
\]
\[
+ (a^2b^2 + b^2c^2 + c^2a^2)(x^2 + y^2 + z^2) - a^2b^2c^2 = 0,
\]
c'est-\`a-dire que l'\'equation est du degr\'e 4 en $x^2$, $y^2$, $z^2$, et du degr\'e 16 en $x$, $y$, $z$.

\textbf{22.} En \'ecrivant l'\'equation sous la forme
\[
U = (a^2x^2 + b^2y^2 + c^2z^2)(x^2 + y^2 + z^2)^2 - \ldots = 0,
\]
on voit que la surface est du quatri\`eme ordre par rapport aux carr\'es des
coordonn\'ees. Cette propri\'et\'e est fondamentale pour la th\'eorie de la surface.

\textbf{23.} Les points coniques de la surface sont donn\'es par les \'equations
\[
x = 0, \quad y = 0, \quad z = 0,
\]
et les valeurs correspondantes de $r$ sont $r = a$, $r = b$, $r = c$. Les points coniques
sont donc situ\'es sur les axes de coordonn\'ees, \`a des distances $a$, $b$, $c$ de l'origine.

\textbf{24.} Les ar\^etes de rebroussement de la surface sont des courbes du huiti\`eme
ordre, qui passent par les points coniques. Ces courbes sont les lieux des points
o\`u le plan tangent est stationnaire.

\textbf{25.} La surface des ondes peut \^etre g\'en\'eralis\'ee en consid\'erant un nombre
quelconque de termes dans l'\'equation fondamentale. On obtient ainsi une classe
de surfaces qui comprennent la surface des ondes comme cas particulier.

\textbf{26.} On peut aussi \'etudier la surface des ondes au moyen de la th\'eorie
des invariants. Les invariants de la surface sont des fonctions des coefficients
$a$, $b$, $c$, qui restent inchang\'ees par les transformations orthogonales.

\textbf{27.} La surface des ondes est un cas particulier de la surface de Kummer,
qui est une surface du quatri\`eme ordre ayant 16 points singuliers. La surface de
Kummer a \'et\'e \'etudi\'ee en d\'etail par son auteur, et ses propri\'et\'es sont fondamentales
pour la g\'eom\'etrie alg\'ebrique.

\textbf{28.} Je m'arr\^ete pour un moment pour consid\'erer la m\^eme th\'eorie par
rapport \`a la surface des centres. La surface des centres est le lieu des centres de
courbure de la surface des ondes; elle est du vingt-quatri\`eme ordre, et ses
propri\'et\'es sont intimement li\'ees \`a celles de la surface des ondes.

\textbf{29.} Il para\^it ainsi que l'ordre de la surface des centres de la surface des
ondes doit \^etre $= 38$. Mais ce r\'esultat doit \^etre v\'erifi\'e par un calcul direct; et il
est possible que l'ordre soit en r\'ealit\'e moindre, \`a cause des singularit\'es de la
surface des ondes.

\textbf{30.} L'\'etude compl\`ete de la surface des ondes et de ses surfaces annexes
pr\'esente encore de nombreux probl\`emes int\'eressants. La m\'ethode employ\'ee dans
cette note peut \^etre appliqu\'ee avec succ\`es \`a plusieurs de ces probl\`emes; et elle
conduit \`a des r\'esultats qui seront utiles pour la th\'eorie de la propagation de la
lumi\`ere dans les cristaux.



\end{document}
